\documentclass[11pt,a4paper]{article}

\usepackage[T1]{fontenc}
\usepackage[utf8]{inputenc}
\usepackage{lmodern}
\usepackage{amsmath,amssymb,mathtools,bm}
\usepackage{booktabs,longtable,tabularx,array}
\usepackage[margin=2.35cm]{geometry}
\usepackage{xcolor}
\usepackage{enumitem}
\usepackage{listings}
\usepackage{microtype}
\usepackage{hyperref}

\hypersetup{
  colorlinks=true,
  linkcolor=blue!55!black,
  citecolor=blue!55!black,
  urlcolor=blue!65!black,
  pdftitle={First-Principles Wick-Contraction Derivations for Pion and Proton Measurements},
  pdfauthor={PyQUDA Measurement physics audit},
}

\allowdisplaybreaks
\setlength{\parindent}{0pt}
\setlength{\parskip}{0.55em}
\setlength{\emergencystretch}{3em}
\setlist[itemize]{topsep=0.3em,itemsep=0.15em}
\setlist[enumerate]{topsep=0.3em,itemsep=0.15em}
\numberwithin{equation}{section}

\definecolor{codegray}{RGB}{245,246,248}
\definecolor{auditred}{RGB}{150,25,25}
\definecolor{auditgreen}{RGB}{20,105,55}
\definecolor{auditamber}{RGB}{145,90,0}
\lstset{
  basicstyle=\ttfamily\small,
  backgroundcolor=\color{codegray},
  frame=single,
  rulecolor=\color{black!15},
  breaklines=true,
  columns=fullflexible,
  keepspaces=true,
  showstringspaces=false,
}

\newcommand{\ii}{\mathrm{i}}
\newcommand{\ee}{\mathrm{e}}
\newcommand{\dd}{\mathrm{d}}
\newcommand{\Tr}{\operatorname{Tr}}
\newcommand{\Trsc}{\operatorname{Tr}_{s,c}}
\newcommand{\Trc}{\operatorname{Tr}_{c}}
\newcommand{\avg}[1]{\left\langle #1\right\rangle}
\newcommand{\abs}[1]{\left|#1\right|}
\newcommand{\norm}[1]{\left\lVert#1\right\rVert}
\newcommand{\cO}{\mathcal{O}}
\newcommand{\cW}{\mathcal{W}}
\newcommand{\cD}{\mathcal{D}}
\newcommand{\cX}{\mathcal{X}}
\newcommand{\Dlr}{\overleftrightarrow D}
\newcommand{\slDlr}{\overleftrightarrow{\not\!D}}
\newcommand{\Gsrc}{\Gamma_{\mathrm{src}}}
\newcommand{\Gsink}{\Gamma_{\mathrm{sink}}}
\newcommand{\Gins}{\Gamma_{\mathrm{ins}}}
\newcommand{\Pf}{\boldsymbol p_f}
\newcommand{\PiMom}{\boldsymbol p_i}
\newcommand{\qv}{\boldsymbol q}
\newcommand{\Dv}{\boldsymbol\Delta}
\newcommand{\Pv}{\boldsymbol P}
\newcommand{\code}[1]{%
  \ifmmode\text{\texttt{\detokenize{#1}}}\else\nolinkurl{#1}\fi}
\newcommand{\good}[1]{\textcolor{auditgreen}{\textbf{#1}}}
\newcommand{\warn}[1]{\textcolor{auditamber}{\textbf{#1}}}
\newcommand{\bad}[1]{\textcolor{auditred}{\textbf{#1}}}
\newcolumntype{L}[1]{>{\raggedright\arraybackslash}p{#1}}

\title{\textbf{First-Principles Wick-Contraction Derivations}\\
  \Large Pion and Proton Measurements in \texttt{pyquda\_measurement\_utils}}
\author{Physics-audit technical note}
\date{Audit baseline:
  \texttt{52927d03f30c8149088f933748e49085ce920b5c}\\
  July 2026}

\begin{document}

\maketitle

\begin{abstract}
This note derives the pion and proton measurement kernels in
\code{pyquda_measurement_utils} from the target continuum matrix elements,
through Euclidean spectral decompositions, flavor-labelled Grassmann
integrals, propagator and sequential-source construction, and finally the
spin--color contractions implemented by the production \code{einsum}
expressions.  It covers two-point functions, electromagnetic form factors
(EMFF), quark and gluon energy--momentum tensor (EMT) building blocks,
qTMD/qGTMD, qPDF/qGPD, pion qTMD wave functions and qDA, stochastic
disconnected loops, the legacy pion soft-factor four-point function, and
background-response measurements.

The central Wick results are that the pion backward line and trace are
correct, the proton direct--exchange determinant and both flavor-specific
sequential kernels are correct, and the qTMDWF/qDA endpoint orientation is
correct.  A physical sign is missing when the saved positive stochastic trace
in the disconnected qTMD/PDF workflow is interpreted as
\(\bar q\Gamma\cW q\).  Most remaining limitations separate a raw Euclidean
correlator from a renormalized and matched physical observable; they are not
additional Wick-contraction failures.
\end{abstract}

\tableofcontents
\newpage

\section{Scope, baseline, and meaning of ``measurement''}

\subsection{Audited source and document purpose}

All formula-to-code statements in this document refer to repository commit
\code{52927d03f30c} (the full hash appears on the title page).  The principal production
files are:
\begin{itemize}
  \item \code{pion_utils_vibe_develop.py},
        \code{proton_utils_vibe_develop.py}, and \code{bw_seq_pyquda.py};
  \item \code{pion_EMFF_vibe_develop.py},
        \code{pion_EMT_vibe_develop.py}, and
        \code{proton_EMT_vibe_develop.py};
  \item \code{pion_qTMD_vibe_develop.py},
        \code{proton_qTMD_pyquda.py},
        \code{pion_qTMDWF_pyquda.py}, and
        \code{qtmd_operator_utils.py};
  \item \code{Disconnected_1pt_qTMD_vibe_develop.py},
        \code{Disconnected_1pt_EMT_vibe_develop.py}, and
        \code{flowed_quark_ringed_norm.py};
  \item \code{pion_soft_factor_vibe_develop.py} and
        \code{pion_current_background_response_vibe_develop.py}.
\end{itemize}

This is a derivation and audit note, not an analysis prescription with a
single universal ratio.  The source and sink interpolators, polarization
projectors, temporal boundary effects, excited states, and the renormalization
scheme must be retained in an ensemble analysis.

\subsection{Three distinct output levels}

The word observable is often used for three different objects:
\begin{enumerate}
  \item a \emph{raw Euclidean correlator}, such as \(C_2\), \(C_3\), or a
        stochastic one-point trace;
  \item a \emph{bare lattice matrix element}, obtained after a spectral fit or
        ratio removes overlaps and Euclidean exponentials;
  \item a \emph{physical observable}, after charge/flavor combination,
        improvement, renormalization, soft subtraction, continuum matching,
        scale evolution, and, for the flowed EMT, the small-flow-time
        expansion.
\end{enumerate}
The audited measurement code predominantly produces level 1.  In correctly
chosen kinematics its output can be used to obtain level 2.  None of EMFF,
EMT, qTMD/qGTMD, qPDF/qGPD, qTMDWF, or qDA is automatically a level-3
observable in the current producer.

\begin{longtable}{L{0.16\textwidth}L{0.31\textwidth}L{0.43\textwidth}}
\toprule
Measurement & Continuum target & Object directly produced \\
\midrule
\endfirsthead
\toprule
Measurement & Continuum target & Object directly produced \\
\midrule
\endhead
Pion EMFF &
\(\avg{\pi(p_f)|J_\mu^{\rm em}|\pi(p_i)}\) and \(F_\pi(Q^2)\) &
One-flavor connected local-current \(C_3\), without the full charge/flavor
combination, \(Z_V\), improvement, disconnected term, or ratio.\\
Proton EMFF &
\(F_1(Q^2)\) and \(F_2(Q^2)\) &
No dedicated producer/extraction; the local limit of the common proton
insertion kernel is only the raw connected building block.\\
Quark EMT &
\(\avg{H(p_f)|T_{\mu\nu}^q|H(p_i)}\) and gravitational form factors &
Connected and disconnected flowed local/one-derivative primitives.\\
Gluon EMT &
\(\avg{H(p_f)|T_{\mu\nu}^g|H(p_i)}\) &
Flowed \(2\Tr_c\sum_\rho F_{\mu\rho}F_{\nu\rho}\), not a complete matched EMT.\\
qTMD/\allowbreak qGTMD &
Staple-bilocal hadron matrix element and its factorized light-cone quantity &
Bare coordinate-space connected or stochastic trace correlator.\\
qPDF/qGPD &
Straight-bilocal quasi distribution/generalized distribution &
Bare coordinate-space connected or stochastic trace correlator.\\
qTMDWF/\allowbreak qDA &
Vacuum-to-pion bilocal wave-function/distribution-amplitude matrix element &
Bare vacuum-to-pion correlator; qTMDWF is currently CG/fixed-gauge.\\
Soft factor &
Soft subtraction compatible with the chosen TMD factorization scheme &
A legacy pion wall-source four-propagator \(C_4\), whose universality for the
GI staple scheme is not established by the contraction alone.\\
Response &
Derivatives of correlators with respect to external operator couplings &
Ordered \(SOS\) and \(SO_2SO_1S\) response propagators for diagnostics.\\
\bottomrule
\end{longtable}

\section{Euclidean, array, Fourier, and spectral conventions}

\subsection{Gamma matrices and propagator orientation}

The code uses the PyQUDA/DeGrand--Rossi Euclidean basis,
\begin{equation}
 \{\gamma_\mu,\gamma_\nu\}=2\delta_{\mu\nu},\qquad
 \gamma_\mu^\dagger=\gamma_\mu,\qquad
 \gamma_5=\gamma_1\gamma_2\gamma_3\gamma_4,
\end{equation}
and the Wilson--clover operator satisfies
\begin{equation}
 D^\dagger=\gamma_5D\gamma_5,\qquad
 S(y,x)=D^{-1}(y,x)=\gamma_5S(x,y)^\dagger\gamma_5.
 \label{eq:g5-hermiticity}
\end{equation}
For one flavor,
\begin{equation}
 S_{\alpha\beta}^{ab}(x,y)
 =\avg{q_\alpha^a(x)\bar q_\beta^b(y)}.
\end{equation}
The production propagator layout is
\begin{lstlisting}
prop.data[w,t,z,y,x_cb, spin_sink, spin_source,
          color_sink, color_source]
\end{lstlisting}
so the final four indices are ordered as
\(S_{\alpha\beta}^{ab}\).  A Hermitian conjugate must exchange \emph{both}
spin and color endpoints.

The common pion backward-line expression is
\begin{lstlisting}
einsum("ij,wtzyxilab,kl->wtzyxkjba",
       gamma5, prop.data.conj(), gamma5)
\end{lstlisting}
and therefore returns
\begin{align}
 B_{\delta\gamma}^{ba}(x)
 &=(\gamma_5)_{\delta\ell}
   \bigl[S_{\rho\ell}^{ab}(x,x_0)\bigr]^*
   (\gamma_5)_{\rho\gamma}
 \nonumber\\
 &=S_{\delta\gamma}^{ba}(x_0,x).
 \label{eq:backward-line-indices}
\end{align}
This is a full compound spin--color dagger, not a color-only transpose.

\subsection{Propagator construction and smearing}

For flavor \(f\),
\begin{equation}
 D_f S_f(\,\cdot\,,x_0)=\eta_{x_0},\qquad
 S_f=D_f^{-1}.
\end{equation}
Writing a spatial smearing operator as \(\mathcal S_{\boldsymbol k}\), the
source- and sink-smeared line has the schematic form
\begin{equation}
 S_f^{XY}(x,x_0)
 =\mathcal S_X(x)\,D_f^{-1}(x,x_0)\,
  \mathcal S_Y^\dagger(x_0).
 \label{eq:smeared-prop}
\end{equation}
Momentum-boosted smearing modifies overlap with a moving state.  It does not
replace the Fourier momentum projection and by itself does not set the
physical hadron momentum.  The positive- and negative-boost pion lines may be
chosen as active or spectator lines, but flavor labels must be assigned before
the isospin-degenerate equality \(S_u=S_d\) is used.

\subsection{Fourier phases and the actual external kinematics}

\code{MomentumPhase} constructs
\begin{equation}
 \Phi_{\boldsymbol n}(x;x_0)
 =\exp\left[
 +2\pi\ii\sum_{j=1}^{3}
 \frac{n_j(x_j-x_{0j})}{L_j}
 \right].
 \label{eq:momentum-phase}
\end{equation}
Two-point applications pass \(-\boldsymbol p\), which gives the standard
\(\exp[-\ii\boldsymbol p\cdot(\boldsymbol x-\boldsymbol x_0)]\) projector.
The meson or baryon fixed-sink sequential source contains
\(\Phi_{\boldsymbol p_f}(y)\).  The subsequent complex conjugation used to
form the backward sequential line turns it into the effective sink projector
\(\exp[-\ii\boldsymbol p_f\cdot(\boldsymbol y-\boldsymbol x_0)]\).
The insertion is multiplied by \(\Phi_{\boldsymbol q}(x)\).

Before projection, translation covariance gives
\begin{equation}
 C_3(y,x;x_0)\sim
 \ee^{+\ii\boldsymbol p_f\cdot(\boldsymbol y-\boldsymbol x)}
 \ee^{+\ii\boldsymbol p_i\cdot(\boldsymbol x-\boldsymbol x_0)}.
\end{equation}
Multiplication by the two code phases leaves a spatial sum only if
\begin{equation}
 \boxed{\boldsymbol q=\boldsymbol p_f-\boldsymbol p_i,\qquad
        \boldsymbol p_i=\boldsymbol p_f-\boldsymbol q.}
 \label{eq:kinematics}
\end{equation}
The module docstrings and dedicated measurement notes use the executed
plus-sign insertion convention in Eq.~\eqref{eq:momentum-phase}.  A sequential
source is built with \(+\boldsymbol p_f\); complex conjugation in the backward
line produces the effective sink projector with \(-\boldsymbol p_f\).  The
numerical kinematics is Eq.~\eqref{eq:kinematics}.

For off-forward bilocals,
\begin{equation}
 \Delta=p_f-p_i=q,\qquad
 P=\frac{p_f+p_i}{2}=p_f-\frac q2,
 \label{eq:P-Delta}
\end{equation}
\begin{equation}
 t=\Delta^2=(E_f-E_i)^2-\abs{\boldsymbol\Delta}^{2},
 \qquad \xi=-\frac{\Delta^+}{2P^+}.
 \label{eq:t-xi}
\end{equation}
Thus a nonzero \(\boldsymbol q\) is qGTMD/qGPD kinematics, not a forward
qTMD/qPDF measurement.

\subsection{Spectral decompositions}

With relativistically normalized finite-volume states,
\begin{equation}
 C_2^H(t,\boldsymbol p)
 =\sum_n\frac{Z_n(\boldsymbol p)Z_n^\dagger(\boldsymbol p)}{2E_n}
 \ee^{-E_nt}
 \text{backward/wrap contributions}.
 \label{eq:C2-spectral}
\end{equation}
For \(0<\tau<t_{\rm sep}\),
\begin{align}
 C_3^H(\tau,t_{\rm sep};p_f,p_i)
 &\longrightarrow
 \frac{Z_f(\boldsymbol p_f)Z_i^\dagger(\boldsymbol p_i)}
      {4E_fE_i}
 \ee^{-E_f(t_{\rm sep}-\tau)}
 \ee^{-E_i\tau}
 \nonumber\\
 &\hspace{2.5cm}\times
 \avg{H(p_f)|\cO|H(p_i)}.
 \label{eq:C3-spectral}
\end{align}
The producer does not perform the spectral fit or construct a complete
form-factor/distribution ratio.  An antiperiodic quark temporal boundary does
not remove hadron wrap-around terms.

For a vacuum-to-pion bilocal,
\begin{align}
 C_\Gamma(t,\boldsymbol P,b)
 &=\sum_{\boldsymbol x}
   \ee^{-\ii\boldsymbol P\cdot(\boldsymbol x-\boldsymbol x_0)}
   \avg{\cO_\Gamma(\boldsymbol x,t;b)
        \cO_\pi^\dagger(x_0)}
 \nonumber\\
 &\longrightarrow
 \frac{\avg{0|\cO_\Gamma(b)|\pi(P)}
       Z_\pi^\dagger(\boldsymbol P)}
      {2E_\pi}
 \ee^{-E_\pi t}.
 \label{eq:vac-pion-spectral}
\end{align}

\section{Grassmann integration and universal loop signs}

\subsection{Generating functional}

For one flavor and sources \(\eta,\bar\eta\),
\begin{align}
 Z[\eta,\bar\eta]
 &=\int DqD\bar q\,
   \exp\left[-\bar qDq+\bar\eta q+\bar q\eta\right]
 \nonumber\\
 &=\det D\,\exp\left[\bar\eta D^{-1}\eta\right].
\end{align}
Differentiating in a fixed source order gives
\begin{equation}
 \avg{q_i\bar q_j}=S_{ij},\qquad
 \avg{q_iq_k\bar q_j\bar q_l}
 =S_{ij}S_{kl}-S_{il}S_{kj}.
 \label{eq:two-identical-wick}
\end{equation}
Equation~\eqref{eq:two-identical-wick} is a determinant over complete endpoint
labels.  If \(i=(\alpha,a)\) is a spin--color compound index, exchange means
exchanging the full pair \((\alpha,a)\).

\subsection{Closed-fermion-loop minus sign}

Introduce a bilinear perturbation:
\begin{equation}
 Z[\lambda]
 =\int DqD\bar q\,
   \ee^{-\bar q(D+\lambda M)q}
 =\det(D+\lambda M).
\end{equation}
Then
\begin{equation}
 \left.\frac{\partial\ln Z}{\partial\lambda}\right|_0
 =\Tr(D^{-1}M)
 =-\avg{\bar qMq},
\end{equation}
and hence
\begin{equation}
 \boxed{\avg{\bar qMq}=-\Tr[MS].}
 \label{eq:closed-loop-sign}
\end{equation}
For noise satisfying
\(\mathbb E[\xi_i\xi_j^*]=\delta_{ij}\), let
\(\eta=D^{-1}\xi\).  The stochastic identity is
\begin{equation}
 \mathbb E_\xi[\xi^\dagger M\eta]
 =\Tr[MD^{-1}].
 \label{eq:stochastic-trace}
\end{equation}
Equation~\eqref{eq:stochastic-trace} estimates a positive matrix trace; it
does not generate the physical closed-loop minus in
Eq.~\eqref{eq:closed-loop-sign}.

\section{Pion two- and three-point Wick contractions}

\subsection{Pion interpolators and complete C2 derivation}

Use a charged-pion sink
\begin{equation}
 \cO_s(x)=\bar d_\gamma^c(x)
 (\Gamma_s)_{\gamma\alpha}u_\alpha^c(x)
\end{equation}
and source
\begin{equation}
 \cO_0^\dagger(x_0)
 =\bar u_\beta^b(x_0)
 (\Gamma_0)_{\beta\delta}d_\delta^b(x_0).
\end{equation}
The field order is
\begin{equation}
 \bar d(x),\ u(x),\ \bar u(x_0),\ d(x_0).
\end{equation}
Flavor conservation permits only
\begin{equation}
 u_\alpha^c(x)\bar u_\beta^b(x_0)
 \longrightarrow
 (S_u)_{\alpha\beta}^{cb}(x,x_0),
\end{equation}
\begin{equation}
 d_\delta^b(x_0)\bar d_\gamma^c(x)
 \longrightarrow
 (S_d)_{\delta\gamma}^{bc}(x_0,x).
\end{equation}
Even if \(m_u=m_d\), pairings such as \(u\bar d\) vanish because the Dirac
operator is flavor diagonal.  Moving the Grassmann variables into paired
order gives the connected-meson sign,
\begin{equation}
 C_{2,\mathrm{formal}}^\pi(x)
 =-\Trsc\left[
 S_d(x_0,x)\Gamma_sS_u(x,x_0)\Gamma_0
 \right].
 \label{eq:pion-C2-formal}
\end{equation}

For the Euclidean pseudoscalar,
\begin{align}
 (\bar d\gamma_5u)^\dagger
 &=\bar u\,\gamma_4\gamma_5^\dagger\gamma_4d
 =-\bar u\gamma_5d.
\end{align}
This adjoint minus cancels the Grassmann minus.  With the production source
convention \(\Gamma_{\rm src}=+\gamma_5\), the implemented correlator is
\begin{equation}
 C_2^\pi(x)
 =\Trsc\left[
 S_d(x_0,x)\gamma_5S_u(x,x_0)\gamma_5
 \right].
 \label{eq:pion-C2-production}
\end{equation}
Using Eq.~\eqref{eq:g5-hermiticity}, equal degenerate lines give
\begin{equation}
 C_2^\pi(x)
 =\Trsc[S(x,x_0)^\dagger S(x,x_0)]\geq0.
 \label{eq:pion-positive-norm}
\end{equation}

\subsection{Index map to the pion production contraction}

Let the backward line returned by Eq.~\eqref{eq:backward-line-indices} be
\(B_{ji}^{cf}\).  The production operations are
\begin{align}
 X_{jm}^{cf}
 &=B_{ji}^{cf}(\Gamma_s)_{im},\\
 C_{\rm site}
 &=X_{ji}^{ab}
   (S_u)_{il}^{ba}
   (\Gamma_{\rm src})_{lj}.
\end{align}
The actual expressions are:
\begin{lstlisting}
sink_inserted =
  einsum("wtzyxjicf,im->wtzyxjmcf",
         backward_line, sink_gamma)
corr_site =
  einsum("wtzyxjiab,wtzyxilba,lj->wtzyx",
         sink_inserted, forward_prop, source_gamma)
\end{lstlisting}
Every spin and color index appears exactly twice, in the Wick order of
Eq.~\eqref{eq:pion-C2-production}.  The final contraction with
\code{"qwtzyx,wtzyx->qt"} applies the Fourier projector and sums the spatial
sites.

\subsection{Connected pion C3 with a bilocal insertion}

For an active-\(u\) operator based at \(x\),
\begin{equation}
 J_\Gamma^u(x;b)
 =\bar u(x)\Gamma\cW[x,x+b]u(x+b),
 \label{eq:pion-bilocal}
\end{equation}
the flavor-preserving connected Wick path is
\begin{align}
 C_{3,u}^\pi
 &=\sum_{\boldsymbol x,\boldsymbol y}
 \ee^{-\ii\boldsymbol p_f\cdot(\boldsymbol y-\boldsymbol x_0)}
 \ee^{+\ii\boldsymbol q\cdot(\boldsymbol x-\boldsymbol x_0)}
 \nonumber\\
 &\quad\times\Trsc\left[
 S_d(x_0,y)\Gsink
 S_u(y,x)\Gamma
 \cW[x,x+b]S_u(x+b,x_0)\Gsrc
 \right].
 \label{eq:pion-C3-master}
\end{align}
The antiquark \(d\) line is the spectator.  An insertion on the \(d\) line
has a separately flavor-labelled path; isospin symmetry may relate the two
only after both paths have been defined.

\subsection{Meson fixed-sink sequential source}

The sink sum in Eq.~\eqref{eq:pion-C3-master} is absorbed into a sequential
solve.  With the spectator source-to-sink propagator, the source placed on
\(y_4=t_0+t_{\rm sep}\) is
\begin{equation}
 \eta_{\rm seq}(y)
 =\delta_{y_4,t_0+t_{\rm sep}}\,
 \Phi_{\boldsymbol p_f}(y;x_0)\,
 \left[\gamma_5\Gsink^\dagger\gamma_5\right]
 S_{\rm spectator}(y,x_0).
 \label{eq:meson-seq-source}
\end{equation}
The code constructs
\begin{equation}
 \Gamma_{\rm seq}
 =\gamma_5\Gsink^\dagger\gamma_5
\end{equation}
with
\begin{lstlisting}
gamma_seq = einsum("ab,cb,cd->ad",
                   gamma5, sink_gamma.conj(), gamma5)
\end{lstlisting}
and solves
\begin{equation}
 D S_{\rm seq}=\eta_{\rm seq}.
\end{equation}
The backward sequential line
\begin{equation}
 B_{\rm seq}(x)=\gamma_5S_{\rm seq}(x)^\dagger\gamma_5
\end{equation}
contains the conjugated sink phase and the summed line
\(S_d(x_0,y)\Gsink S_u(y,x)\).  The remaining local trace is
\begin{equation}
 \Trsc\left[
 B_{\rm seq}(x)\Gamma\cW[x,x+b]
 S_{\rm active}(x+b,x_0)\Gsrc
 \right],
 \label{eq:pion-C3-after-seq}
\end{equation}
which is precisely Eq.~\eqref{eq:pion-C3-master}.

\section{Proton Wick determinant and sequential cuts}

\subsection{Interpolating field}

Write the proton sink as
\begin{equation}
 \chi_m(x)
 =\epsilon_{abc}G_{ij}
 u_i^a(x)d_j^b(x)u_m^c(x),
 \qquad G=C\Gamma_{\rm interp},
 \label{eq:proton-interpolator}
\end{equation}
with \(C=\ii\gamma_2\gamma_4\).  Production supports
\(C\gamma_5\), \(C\gamma_4\gamma_5\), and
\(C\gamma_3\gamma_5\), labelled \code{5}, \code{T5}, and \code{Z5}.
All supported \(G\) matrices are antisymmetric in the adopted basis,
\begin{equation}
 G^T=-G.
\end{equation}
The projected source contains
\begin{equation}
 \bar\chi_n(x_0)\ \propto\
 \epsilon_{def}G_{kl}
 \bar u_k^d(x_0)\bar d_l^e(x_0)\bar u_n^f(x_0),
\end{equation}
where the fixed Euclidean-conjugation convention is absorbed in \(G\) and the
external spin projector \(P_{mn}\).  For positive-parity unpolarized
projection, \(P_+=(1+\gamma_4)/4\).

\subsection{Flavor-labelled Wick expansion}

The \(d\) flavor has the unique pairing
\begin{equation}
 d_j^b(x)\bar d_l^e(x_0)
 \longrightarrow
 (S_d)_{jl}^{be}.
\end{equation}
The two identical \(u\) fields give the determinant
\begin{equation}
 \det\begin{pmatrix}
 (S_u)_{ik}^{ad} & (S_u)_{in}^{af}\\
 (S_u)_{mk}^{cd} & (S_u)_{mn}^{cf}
 \end{pmatrix}
 =
 (S_u)_{ik}^{ad}(S_u)_{mn}^{cf}
 -(S_u)_{in}^{af}(S_u)_{mk}^{cd}.
 \label{eq:proton-u-determinant}
\end{equation}
The exchange in the second product swaps complete spin--color source
endpoints.  Including the common field-order convention used by production,
\begin{equation}
 \boxed{\mathcal C_p[S_u,S_d]=-\mathcal D+\mathcal X,}
 \label{eq:proton-C2-master}
\end{equation}
where
\begin{align}
 \mathcal D
 &=\epsilon_{abc}\epsilon_{def}
   G_{ij}G_{kl}P_{mn}
   (S_u)_{ik}^{ad}(S_d)_{jl}^{be}(S_u)_{mn}^{cf},
 \label{eq:proton-direct}\\
 \mathcal X
 &=\epsilon_{abc}\epsilon_{def}
   G_{ij}G_{kl}P_{mn}
   (S_u)_{in}^{af}(S_d)_{jl}^{be}(S_u)_{mk}^{cd}.
 \label{eq:proton-exchange}
\end{align}
The equality \(S_u=S_d\) may be imposed for degenerate masses after
Eqs.~\eqref{eq:proton-direct}--\eqref{eq:proton-exchange} are derived.  It
does not create an \(u\)-to-\(\bar d\) contraction.

\subsection{Why the printed second term is the complete exchange}

The scalar represented by the second production \code{einsum} is
\begin{equation}
 T_2^{\rm code}
 =\epsilon_{abc}\epsilon_{def}
 G_{ij}G_{kl}P_{mn}
 S_{ik}^{ad}S_{jn}^{be}S_{ml}^{cf}.
 \label{eq:proton-printed-term2}
\end{equation}
The middle propagator in Eq.~\eqref{eq:proton-printed-term2} can look like a
partial \(u/d\) exchange if the dummy labels are interpreted as persistent
flavor labels.  They are not.  For a single degenerate propagator tensor,
simultaneous relabelling of source spin and color dummy indices through both
epsilon tensors, followed by \(G^T=-G\), yields
\begin{equation}
 T_2^{\rm code}=-\mathcal X\qquad(S_u=S_d=S).
 \label{eq:term2-exchange-equivalence}
\end{equation}
Consequently the production combination
\begin{lstlisting}
corr = -term1 - term2
\end{lstlisting}
is
\begin{equation}
 -\mathcal D-(-\mathcal X)=-\mathcal D+\mathcal X.
\end{equation}
An independent explicit compound-index sum found
\begin{align}
 \abs{T_2^{\rm code}+\mathcal X}
 &=5.66\times10^{-14},\\
 \abs{C_{\rm code}-(-\mathcal D+\mathcal X)}
 &=5.55\times10^{-14}
\end{align}
for random complex propagators.  A spin-only or
color-only swap does not satisfy this identity.

\subsection{Flavor-specific sequential kernels as derivatives}

The mathematically unambiguous definition of a cut baryon line is the
holomorphic derivative of the same flavor-labelled polynomial:
\begin{equation}
 (K_d)_{\alpha\beta}^{ab}
 =\frac{\partial\mathcal C_p[S_u,S_d]}
        {\partial(S_d)_{\alpha\beta}^{ab}},
 \qquad
 (K_u)_{\alpha\beta}^{ab}
 =\frac{\partial\mathcal C_p[S_u,S_d]}
        {\partial(S_u)_{\alpha\beta}^{ab}}.
 \label{eq:sequential-kernel-derivative}
\end{equation}
Because \(\mathcal C_p\) is linear in \(S_d\), \(K_d\) has two terms: the
derivative of the direct and exchange products.  Because it is quadratic in
\(S_u\), \(K_u\) has four terms: both \(u\) factors can be cut in each of the
direct and exchange products.  This explains the two terms in
\code{down_quark_insertion_pyquda} and the four terms in
\code{up_quark_insertion_pyquda}; the latter already includes the
multiplicity of the two identical \(u\) lines.

The line converted into a Dirac source has reversed compound endpoints:
\begin{equation}
 K_{\rm source}[\beta,\alpha,b,a]
 =\frac{\partial\mathcal C_p}
        {\partial S[\alpha,\beta,a,b]}.
 \label{eq:kernel-source-transpose}
\end{equation}
Both the spin and color pairs are transposed.  Exact polynomial
differentiation gives maximum discrepancies
\begin{equation}
 \norm{K_d^{\rm code}-K_d^{\rm exact}}_\infty
 =1.99\times10^{-15},\qquad
 \norm{K_u^{\rm code}-K_u^{\rm exact}}_\infty
 =2.08\times10^{-14}.
\end{equation}
The production-adapter comparison gives approximately
\(3.6\times10^{-15}\) for each flavor.

\subsection{Proton C3 after the sequential solve}

After inverting the flavor-specific sequential source, a local or displaced
insertion has the common form
\begin{equation}
 C_{3,f}
 =\sum_{\boldsymbol x}
 \ee^{+\ii\boldsymbol q\cdot(\boldsymbol x-\boldsymbol x_0)}
 \Trsc\left[
 B_{{\rm seq},f}(x)\Gamma
 \cW[x,x+b]S_f(x+b,x_0)
 \right].
 \label{eq:proton-C3-after-seq}
\end{equation}
The connected nucleon runner uses
\begin{lstlisting}
einsum("pwtzyxjicf,gim,wtzyxmjfc->pgwtzyx",
       sequential, gamma_ls, current_prop)
\end{lstlisting}
which maps exactly to
\begin{equation}
 (B_{{\rm seq},f})_{ji}^{cf}\,
 \Gamma_{im}\,
 (\cW S_f)_{mj}^{fc}.
\end{equation}
The same cut flavor closes both spin and color.  No additional factor two
may be applied to the up channel in downstream analysis.

\section{Electromagnetic form factors}

\subsection{Continuum matrix elements}

The electromagnetic current is
\begin{equation}
 J_\mu^{\rm em}
 =\sum_f e_f\,\bar q_f\gamma_\mu q_f.
 \label{eq:em-current}
\end{equation}
For a charged pion,
\begin{equation}
 \avg{\pi(p_f)|J_\mu^{\rm em}|\pi(p_i)}
 =(p_f+p_i)_\mu F_\pi(Q^2),\qquad
 Q^2=-q^2.
 \label{eq:pion-EMFF}
\end{equation}
For the proton,
\begin{equation}
 \avg{p(p_f,s_f)|J_\mu^{\rm em}|p(p_i,s_i)}
 =
 \bar u(p_f,s_f)
 \left[
 \gamma_\mu F_1(Q^2)
 +\frac{\ii\sigma_{\mu\nu}q_\nu}{2M_N}F_2(Q^2)
 \right]
 u(p_i,s_i).
 \label{eq:proton-EMFF}
\end{equation}
Solving Eq.~\eqref{eq:proton-EMFF} for \(F_1,F_2\) requires several current
directions, external momenta, and spin projectors.  The Sachs form factors are
\begin{equation}
 G_E=F_1-\frac{Q^2}{4M_N^2}F_2,\qquad
 G_M=F_1+F_2.
\end{equation}

For Wilson--clover fermions, a local current generally enters as
\begin{equation}
 J_{\mu,R}^f
 =Z_V^f(1+b_Vam_f)
 \left[
 \bar q_f\gamma_\mu q_f
 +ac_V\partial_\nu(\bar q_f\sigma_{\mu\nu}q_f)
 \right],
 \label{eq:improved-vector-current}
\end{equation}
with scheme/action-dependent coefficients.  Neither the flavor charges in
Eq.~\eqref{eq:em-current} nor the factors in
Eq.~\eqref{eq:improved-vector-current} can be inferred from a gamma insertion.

\subsection{Pion local-current contraction}

Setting \(b=0\), \(\cW=1\), and \(\Gamma=\gamma_\mu\) in
Eq.~\eqref{eq:pion-C3-master} gives
\begin{align}
 C_{3,u}^{\pi,\mu}
 &=\sum_{\boldsymbol x,\boldsymbol y}
 \ee^{-\ii\boldsymbol p_f\cdot(\boldsymbol y-\boldsymbol x_0)}
 \ee^{+\ii\boldsymbol q\cdot(\boldsymbol x-\boldsymbol x_0)}
 \nonumber\\
 &\quad\times
 \Trsc\left[
 S_d(x_0,y)\Gsink
 S_u(y,x)\gamma_\mu S_u(x,x_0)\Gsrc
 \right].
 \label{eq:pion-emff-C3}
\end{align}
\code{pion_EMFF_vibe_develop.py} computes
Eq.~\eqref{eq:pion-emff-C3} through the common 16-gamma scan and the meson
sequential line.  The phase relation is
\(\boldsymbol p_i=\boldsymbol p_f-\boldsymbol q\).

This quantity is a one-active-line raw connected \(C_3\), not
\(F_\pi(Q^2)\).  A complete analysis must:
\begin{enumerate}
  \item include \(C_2(\boldsymbol p_f)\) and
        \(C_2(\boldsymbol p_i)\) with matching interpolators/smearing;
  \item cancel overlap factors and exponentials or perform a simultaneous
        spectral fit;
  \item combine the active-\(u\) and active-\(d\) contractions with \(e_u\)
        and \(e_d\), and add any required disconnected contribution;
  \item apply current improvement and \(Z_V\);
  \item divide by the kinematic factor in Eq.~\eqref{eq:pion-EMFF}.
\end{enumerate}
The production application sets its two-point momentum list from
\code{qext}.  For nonzero \(\boldsymbol p_f\), that list need not contain both
\(\boldsymbol p_f\) and
\(\boldsymbol p_f-\boldsymbol q\).  Existing \(C_3\) can remain usable, but
the missing matching \(C_2\) must then be calculated separately.

\subsection{Proton local-current status}

The proton local-current connected building block is obtained from
Eq.~\eqref{eq:proton-C3-after-seq} at \(b=0\) with
\(\Gamma=\gamma_\mu\), separately for \(f=u,d\):
\begin{equation}
 C_{3,\mu}^{p,\rm conn}
 =e_u C_{3,u}^{\mu}+e_d C_{3,d}^{\mu}.
\end{equation}
The common direct--exchange and sequential kernels are suitable for this
construction.  At the audited commit there is, however, no dedicated proton
EMFF producer that defines the complete local-current charge combination,
projector set, kinematic linear system, and \(F_1/F_2\) extraction.  A proton
qTMD file at \(b=0\) must not automatically be labelled \(F_1\) or \(F_2\).

\section{Quark and gluon energy--momentum tensor measurements}

\subsection{Continuum operators and hadron decompositions}

The symmetric quark one-derivative operator used as the lattice target is
\begin{equation}
 \cO_{\mu\nu}^q
 =\frac14\bar q
 \left[
 \gamma_\mu\Dlr_\nu+\gamma_\nu\Dlr_\mu
 \right]q,
 \qquad
 \Dlr_\nu=\overrightarrow D_\nu-\overleftarrow D_\nu.
 \label{eq:quark-EMT}
\end{equation}
Flavor singlet quark operators mix with gluon operators under
renormalization.

For a spin-zero pion, one common decomposition is
\begin{align}
 \avg{\pi(p_f)|T_{\mu\nu}|\pi(p_i)}
 &=
 2P_\mu P_\nu A_\pi(t)
 +\frac12\left(
 \Delta_\mu\Delta_\nu-\delta_{\mu\nu}\Delta^2
 \right)D_\pi(t)
 \nonumber\\
 &\quad
 +2m_\pi^2\delta_{\mu\nu}\bar C_\pi(t),
 \label{eq:pion-EMT-decomposition}
\end{align}
subject to Euclidean continuation and operator-basis conventions.  For a
spin-\(\tfrac12\) proton, a standard form is
\begin{align}
 \avg{p_f|T_{\mu\nu}|p_i}
 &=\bar u(p_f)
 \Bigg[
 A(t)\gamma_{\{\mu}P_{\nu\}}
 +B(t)\frac{\ii P_{\{\mu}\sigma_{\nu\}\rho}\Delta_\rho}{2M_N}
 \nonumber\\
 &\qquad
 +D(t)\frac{\Delta_\mu\Delta_\nu-\delta_{\mu\nu}\Delta^2}{4M_N}
 +\bar C(t)M_N\delta_{\mu\nu}
 \Bigg]u(p_i).
 \label{eq:proton-EMT-decomposition}
\end{align}
Braces denote symmetrization with the convention chosen in the analysis.
The measurement code does not solve either gravitational-form-factor system.

\subsection{Central gauge-covariant derivative}

The lattice shift convention is
\begin{equation}
 (\mathcal T_{+\mu}q)(x)=U_\mu(x)q(x+\hat\mu),\qquad
 (\mathcal T_{-\mu}q)(x)
 =U_\mu^\dagger(x-\hat\mu)q(x-\hat\mu).
 \label{eq:covariant-shifts}
\end{equation}
The dimensionless central derivative implemented by
\code{_covdev_sym_prop} is
\begin{equation}
 aD_\mu q(x)
 =\frac12\left[
 U_\mu(x)q(x+\hat\mu)
 -U_\mu^\dagger(x-\hat\mu)q(x-\hat\mu)
 \right].
 \label{eq:central-covdev}
\end{equation}
Let \(B\) denote the sequential/backward line and \(S\) the active line.  For
a nonsymmetrized gamma--derivative pair, production evaluates
\begin{align}
 \mathcal B_{\mu\nu}
 &=
 \frac12 B\gamma_\mu(D_\nu S)
 -\frac12 (B\overleftarrow D_\nu)\gamma_\mu S.
 \label{eq:emt-right-left}
\end{align}
It then forms
\begin{equation}
 \mathcal O_{\mu\nu}^{\rm code}
 =\frac12\left(\mathcal B_{\mu\nu}
                    +\mathcal B_{\nu\mu}\right).
 \label{eq:emt-symmetrize}
\end{equation}
The two halves in Eq.~\eqref{eq:emt-right-left} and the final half in
Eq.~\eqref{eq:emt-symmetrize} reproduce the \(1/4\) in
Eq.~\eqref{eq:quark-EMT}, in lattice units.  No \(1/a\) is inserted.  There
is no automatic spacetime trace subtraction.

\subsection{Connected pion EMT contraction}

At flow time \(t_f\), denote the flowed active and sequential lines by
\(S(t_f)\) and \(B(t_f)\).  The primitive local field is
\begin{equation}
 L_\Gamma(x;t_f)
 =\Trsc\left[B(x;t_f)\Gamma S(x;t_f)\Gsrc\right].
 \label{eq:pion-emt-local}
\end{equation}
The derivative primitive is
\begin{align}
 L_{\Gamma,\nu}(x;t_f)
 &=\frac12\Trsc\left[
 B\Gamma(D_\nu S)\Gsrc\right]
 -\frac12\Trsc\left[
 (B\overleftarrow D_\nu)\Gamma S\Gsrc\right].
 \label{eq:pion-emt-derivative}
\end{align}
The \code{opt_einsum} expression for the local part,
\begin{lstlisting}
"wtzyxabij,gbn,wtzyxncji,ca->gwtzyx"
\end{lstlisting}
closes backward color/spin \((a,b,i,j)\), gamma \((b,n)\), forward line
\((n,c,j,i)\), and source gamma \((c,a)\) in the order of
Eq.~\eqref{eq:pion-emt-local}.  Fourier projection then supplies the
\(+\boldsymbol q\) insertion phase.  The same contraction with \(D_\nu S\)
and \(B\overleftarrow D_\nu\) gives
Eq.~\eqref{eq:pion-emt-derivative}.

\subsection{Connected proton EMT contraction}

For the proton, the raw flavor-specific sequential propagator already
contains \(K_u\) or \(K_d\).  The local primitive can be written
\begin{equation}
 L_{\Gamma,f}^p(x;t_f)
 =\Trsc\left[B_{{\rm seq},f}(x;t_f)
 \Gamma S_f(x;t_f)\right],
\end{equation}
with the identical right/left derivative construction of
Eq.~\eqref{eq:pion-emt-derivative}.  In
\code{proton_EMT_vibe_develop.py}, \code{_raw_seq_gamma_scalar_fields}
contracts the complex-conjugated raw sequential line with the active
propagator and then with \(\gamma_5\Gamma\).  This is the raw-sequential
representation of the same
\(\gamma_5S_{\rm seq}^\dagger\gamma_5\Gamma S\) trace.  The up/down Wick
multiplicities remain those of Eq.~\eqref{eq:sequential-kernel-derivative}.

\subsection{Stochastic disconnected quark EMT}

The local stochastic primitive is
\begin{equation}
 T_\Gamma(x)
 =\xi^\dagger(x)\Gamma\eta(x),
 \qquad \eta=D^{-1}\xi,
\end{equation}
which estimates \(+\Tr[\Gamma S]\).  For the derivative operator, production
computes a right-acting central difference and reconstructs the left-acting
piece with gamma-five Hermiticity and momentum reversal.  If
\(\Gamma^\sharp=\gamma_5\Gamma^\dagger\gamma_5
=s_\Gamma\Gamma_{\rm partner}\), the stored derivative is
\begin{equation}
 L_{\Gamma,\nu}^{\rm disc}(q)
 =-\frac14\left[
 A_R(\Gamma,\nu;q)
 -s_\Gamma A_R(\Gamma_{\rm partner},\nu;-q)^*
 \right].
 \label{eq:disc-emt-derivative}
\end{equation}
The factor in Eq.~\eqref{eq:disc-emt-derivative} combines the central
derivative, the two-sided \(\Dlr\) normalization, and the physical
closed-fermion-loop minus.  Thus the derived disconnected EMT derivative has
the loop sign included.  By contrast, the separately stored primitive local
bilinear remains the positive trace estimator and must not be interpreted as
\(\avg{\bar q\Gamma q}\) without Eq.~\eqref{eq:closed-loop-sign}.

\subsection{Gradient flow and ringed normalization}

At positive flow time, gauge and fermion fields obey schematically
\begin{equation}
 \partial_t B_\mu=D_\nu G_{\nu\mu},\qquad
 \partial_t\chi=\Delta\chi,\qquad
 \partial_t\bar\chi=\bar\chi\overleftarrow\Delta,
 \label{eq:gradient-flow}
\end{equation}
with \(B_\mu(0)=A_\mu\), \(\chi(0)=q\), and \(\bar\chi(0)=\bar q\).  Flowed
local products are finite building blocks, but flowed fermions require the
appropriate ringed-field normalization, and a physical EMT is obtained by a
small-flow-time expansion.  Schematically,
\begin{equation}
 T_{\mu\nu}^{R}(x)
 =\lim_{t\to0}\left[
 c_g(t)\,\widetilde{\cO}_{\mu\nu}^g(t,x)
 +\sum_f c_f(t)\,\widetilde{\cO}_{\mu\nu}^{q_f}(t,x)
 +\text{trace and mixing operators}
 \right].
 \label{eq:small-flow-time}
\end{equation}
The standalone ringed helper estimates the kinetic normalization input
\begin{equation}
 K(t)
 =\avg{\bar\chi(t,x)\slDlr\chi(t,x)}
 \end{equation}
through the zero-momentum diagonal derivative sum.  In the production
normalization convention,
\begin{equation}
 K_{\rm pervec}(t,\tau)
 =-\frac{2}{V_3}
 \sum_{\mu=1}^{4}
 L_D[\gamma_\mu,\mu,q=0;t,\tau].
 \label{eq:ringed-kinetic}
\end{equation}
This is an input to a ringed factor, not the factor itself, and is not
automatically multiplied into connected or disconnected EMT output.

\subsection{Flowed gluon primitive}

The clover field strength is constructed from four oriented plaquettes:
\begin{equation}
 F_{\mu\nu}^{\rm clover}
 \propto -\ii\left[
 Q_{\mu\nu}-Q_{\mu\nu}^\dagger
 -\frac1{N_c}\Tr_c(Q_{\mu\nu}-Q_{\mu\nu}^\dagger)\mathbf1
 \right].
\end{equation}
The code then measures
\begin{equation}
 G_{\mu\nu}^{\rm code}(t,x)
 =2\Tr_c\sum_{\rho}F_{\mu\rho}(t,x)F_{\nu\rho}(t,x).
 \label{eq:gluon-primitive}
\end{equation}
Equation~\eqref{eq:gluon-primitive} is a useful flowed primitive.  It lacks
the \(\delta_{\mu\nu}F_{\rho\sigma}F_{\rho\sigma}\) trace structure, the
physical Euclidean sign/coefficient convention, and the matching/mixing
coefficients required for a complete gluon EMT.

\subsection{EMT audit interpretation}

The right/left derivative, symmetrization, and lattice \(1/4\) normalization
are implemented consistently for the connected pion/proton and derived
disconnected quark EMT.  The output is nevertheless
\emph{unringed/unrenormalized flowed data}.  Trace combinations, the ringed
factor, small-flow matching, quark--gluon mixing, and a controlled window
\begin{equation}
 a^2\ll 8t\ll\Lambda_{\rm QCD}^{-2}
\end{equation}
belong to the subsequent analysis.

\section{qTMD/qGTMD and qPDF/qGPD}

\subsection{Hadron bilocal matrix element}

A convenient translation-covariant definition of the bare coordinate-space
matrix element is
\begin{equation}
 \widetilde W_\Gamma^f(b;P,\Delta,\eta)
 =
 \avg{H(p_f)|
 \bar q_f(-b/2)\Gamma
 \cW_\eta[-b/2,b/2]q_f(b/2)
 |H(p_i)}.
 \label{eq:generic-bilocal}
\end{equation}
The code uses an equivalent basepoint/end-point form,
\begin{equation}
 \cO_\Gamma^f(x;b)
 =\bar q_f(x)\Gamma\cW[x,x+b]q_f(x+b).
 \label{eq:basepoint-bilocal}
\end{equation}
Translation between Eqs.~\eqref{eq:generic-bilocal} and
\eqref{eq:basepoint-bilocal} only changes the known phase associated with the
operator center; it does not change the Wilson path.

For \(b=b_z\hat z+b_T\hat e_T\) and a staple of extent \(\eta\), the forward
case \(\Delta=0\) is a qTMD coordinate-space matrix element.  When
\(\Delta\neq0\), it is a qGTMD matrix element.  The straight-link limit
\(b_T=0\) is the qPDF matrix element for \(\Delta=0\), and the qGPD matrix
element for \(\Delta\neq0\).

\subsection{Straight Wilson line and covariant transport}

For positive \(b_z\),
\begin{equation}
 W[x,x+b_z\hat z]
 =U_z(x)U_z(x+\hat z)\cdots
 U_z(x+(b_z-1)\hat z).
 \label{eq:straight-positive}
\end{equation}
For negative displacement, the reversed conjugate link product is required:
\begin{equation}
 W[x,x-b\hat z]
 =U_z^\dagger(x-\hat z)
 U_z^\dagger(x-2\hat z)\cdots
 U_z^\dagger(x-b\hat z),\qquad b>0.
 \label{eq:straight-negative}
\end{equation}
Repeated application of Eq.~\eqref{eq:covariant-shifts} produces
\begin{equation}
 (\mathcal T_{+z}^{\,b_z}q)(x)
 =W[x,x+b_z\hat z]q(x+b_z\hat z)
\end{equation}
for positive \(b_z\), and Eq.~\eqref{eq:straight-negative} for negative
\(b_z\).  The qPDF implementation restarts the negative branch from the
undisplaced field at \(b_z=-1\), so every incremental step is nearest
neighbor.

\subsection{Fixed-length GI staple}

The GI qTMD index is
\begin{equation}
 [b_T,b_z,\eta,\mathrm{dir}_T].
\end{equation}
For even \(b_z\), nonnegative \(b_T\), and
\(\eta\geq\abs{b_z}/2\), production decomposes the path into signed segments
\begin{equation}
 \left[
 (z,\eta+b_z/2),\
 (e_T,b_T),\
 (z,b_z/2-\eta)
 \right].
 \label{eq:staple-segments}
\end{equation}
The net endpoint displacement is
\begin{equation}
 (\eta+b_z/2)\hat z+b_T\hat e_T
 +(b_z/2-\eta)\hat z
 =b_z\hat z+b_T\hat e_T,
\end{equation}
and, for the allowed range,
\begin{equation}
 L_{\rm path}
 =\abs{\eta+b_z/2}+b_T+\abs{b_z/2-\eta}
 =2\eta+b_T.
\end{equation}
The implementation applies the segment list in reverse order.  This follows
operator composition:
\begin{equation}
 \mathcal T_{s_1}\mathcal T_{s_2}\mathcal T_{s_3}q
 \quad\hbox{is evaluated by applying }\mathcal T_{s_3}
 \hbox{ first}.
\end{equation}
Because every \(\mathcal T\) pulls an endpoint field back to the base point,
the reverse execution order is the correct physical staple, not a reversed
Wilson line.

At \(b_T=0\) and \(\eta=\abs{b_z}/2\),
Eq.~\eqref{eq:staple-segments} collapses to the positive or negative straight
link.  Odd \(b_z\) is excluded by this half-integer-free fixed-length
parameterization, not by a fundamental restriction on Wilson lines.

\subsection{Gauge covariance and CG versus GI}

Under a local gauge transformation,
\begin{equation}
 q(x)\to\Omega(x)q(x),\quad
 \bar q(x)\to\bar q(x)\Omega^\dagger(x),\quad
 U_\mu(x)\to\Omega(x)U_\mu(x)\Omega^\dagger(x+\hat\mu).
\end{equation}
Every internal \(\Omega^\dagger\Omega\) cancels in a link product:
\begin{equation}
 W[x,y]\to\Omega(x)W[x,y]\Omega^\dagger(y).
\end{equation}
Therefore
\begin{equation}
 \bar q(x)\Gamma W[x,y]q(y)
\to
\bar q(x)\Omega^\dagger(x)\Gamma\Omega(x)
W[x,y]\Omega^\dagger(y)\Omega(y)q(y),
\end{equation}
which is invariant because color transformations commute with spin
\(\Gamma\).  Random-SU(3) tests of the production straight and staple
transporters agree with this proof to \(3.38\times10^{-15}\).

The \code{GI_*} paths use these covariant transports.  The \code{CG_*} paths
use ordinary coordinate shifts and omit \(W\).  They are meaningful only
after a specified gauge fixing, such as Coulomb gauge.  A filename containing
\code{cg} is not a mathematical enforcement of gauge fixing; the method,
tolerance, and convention must be provenance metadata.

\subsection{Connected pion qTMD/qGTMD and qPDF/qGPD}

The pion connected nonlocal correlator is the master expression
Eq.~\eqref{eq:pion-C3-master}.  The shifted active propagator is
\begin{equation}
 S_u^{(b)}(x,x_0)
 =W[x,x+b]S_u(x+b,x_0).
 \label{eq:shifted-active-prop}
\end{equation}
Substitution into Eq.~\eqref{eq:pion-C3-after-seq} yields
\begin{equation}
 C_{3,u}^{\pi,\Gamma}(b)
 =\sum_x\Phi_{\boldsymbol q}(x)
 \Trsc\left[
 B_{\rm seq}(x)\Gamma S_u^{(b)}(x,x_0)\Gsrc
 \right].
 \label{eq:pion-qtmd-code-form}
\end{equation}
\code{contract_qTMD_GI} builds a cached transporter for every staple and
applies Eq.~\eqref{eq:shifted-active-prop}; \code{contract_PDF} incrementally
applies the straight shifts.  Both then call the common pion gamma scan, so
the spin--color Wick closure is identical to the verified pion \(C_2/C_3\)
kernel.

\subsection{Connected proton qTMD/qGTMD and qPDF/qGPD}

For each active flavor, the proton connected expression is
Eq.~\eqref{eq:proton-C3-after-seq} with the same shifted line
Eq.~\eqref{eq:shifted-active-prop}.  The down and up sequential sources are
respectively the exact two-term and four-term cuts of
Eq.~\eqref{eq:proton-C2-master}.  The production contraction
\begin{equation}
 (B_{{\rm seq},f})_{ji}^{cf}
 \Gamma_{im}
 (S_f^{(b)})_{mj}^{fc}
\end{equation}
closes the correct active flavor.  The Wilson geometry and the baryon Wick
determinant are logically separate: correctness of one does not compensate
for an error in the other.  Both pass independent checks here.

\subsection{Forward versus off-forward production kinematics}

From Eqs.~\eqref{eq:kinematics}--\eqref{eq:t-xi}:
\begin{equation}
 \boldsymbol q=0
 \Rightarrow p_i=p_f,\ \Delta=0
 \quad\hbox{(qTMD or qPDF)},
\end{equation}
\begin{equation}
 \boldsymbol q\neq0
 \Rightarrow p_i=p_f-q,\ \Delta=q
 \quad\hbox{(qGTMD or qGPD)}.
\end{equation}
At the audited commit, the pion and proton TMD production runners hard-code
\(\boldsymbol p_f=0\).  Consequently:
\begin{equation}
 p_i=-q,\qquad P=-q/2.
\end{equation}
For the forward list, \(q=0\) also gives \(P_z=0\).  It is a valid raw
zero-momentum bilocal correlator but not the large-\(P_z\) quasi-distribution
regime required for LaMET matching.  Momentum-boosted smearing cannot change
this Fourier-projected external momentum.  Nonzero-\(q\) output is a low- or
zero-average-momentum off-forward matrix element and should be labelled
qGTMD/qGPD with \(P,\Delta,\xi,t\) recorded.

\subsection{From a bare bilocal to a distribution}

The measured coordinate-space object requires, schematically,
\begin{equation}
 \widetilde W_\Gamma^{R}
 =Z_\Gamma^{-1}(b,\eta,a,\mu,\text{scheme})
 \widetilde W_\Gamma^{\rm bare},
\end{equation}
followed by the relevant Fourier transform and quasi-to-light-cone matching.
For qTMD/qGTMD, a compatible soft subtraction and rapidity prescription are
also required:
\begin{equation}
 \widetilde F^{\rm sub}
 \sim
 \frac{\widetilde W^{R}}
      {\sqrt{S^{R}(b_T,\eta)}} ,
\qquad
 \frac{\partial\ln\widetilde F}{\partial\ln\sqrt\zeta}
 =K(b_T,\mu).
\end{equation}
The soft function, rapidity scale \(\zeta\), Collins--Soper kernel \(K\), and
LaMET matching kernel are not harmless overall constants.  The current HDF5
correlators are bare inputs to such an analysis, not physical TMDs, GTMDs,
PDFs, or GPDs.

\section{Disconnected qTMD/PDF loop and vacuum subtraction}

\subsection{Operator trace produced by the code}

For the same straight or staple operator, production solves
\(\eta=D^{-1}\xi\), transports \(\eta\) to
\(\cW\eta\), and contracts
\begin{equation}
 T_\Gamma(b,q)
 =\sum_x\Phi_{\boldsymbol q}(x)\,
 \xi^\dagger(x)\Gamma\cW[x,x+b]\eta(x+b).
 \label{eq:disc-positive-trace}
\end{equation}
Its noise expectation is
\begin{equation}
 \mathbb E_\xi[T_\Gamma]
 =+\Tr\left[
 \Phi_{\boldsymbol q}\Gamma\cW_bD^{-1}
 \right].
 \label{eq:disc-trace-target}
\end{equation}
This agrees with the stored metadata:
\begin{lstlisting}
loop_convention = xi_dagger_Gamma_O_b_eta
trace_target    = Tr[P_qtau Gamma O_b Dinv]
\end{lstlisting}

\subsection{Physical Wick loop}

The physical bilinear loop follows from
Eq.~\eqref{eq:closed-loop-sign}:
\begin{equation}
 L_{\rm phys}(b,q)
 =-\Tr\left[
 \Phi_{\boldsymbol q}\Gamma\cW_bD^{-1}
 \right]
 =-\mathbb E_\xi[T_\Gamma].
 \label{eq:disc-physical-loop}
\end{equation}
For a hadron two-point measurement \(C_2^H\), the disconnected three-point
contribution is the vacuum-subtracted covariance
\begin{align}
 C_{3,\rm disc}^{\rm phys}
 &=
 \avg{C_2^H L_{\rm phys}}
 -\avg{C_2^H}\avg{L_{\rm phys}}
 \nonumber\\
 &=-\left[
 \avg{C_2^H T_\Gamma}
 -\avg{C_2^H}\avg{T_\Gamma}
 \right].
 \label{eq:disc-covariance}
\end{align}
Vacuum subtraction does not remove or alter the global minus.

\subsection{Confirmed sign issue and data consequence}

The qTMD/PDF loop producer correctly stores the positive trace estimator in
Eq.~\eqref{eq:disc-trace-target}.  The accompanying physical-combination
description at the audited commit uses the positive covariance without the
outer minus in Eq.~\eqref{eq:disc-covariance}.  No Wilson-path or gamma
convention supplies a compensating sign.  This is a \bad{confirmed P0
physical-sign issue}.

For current-schema raw data with explicit convention metadata, the correction
is exact:
\begin{equation}
 \code{physical_loop}=-\code{loop_pervec}.
\end{equation}
Inversions need not be repeated; any already constructed positive-trace
covariance must be multiplied by \(-1\) and recomputed with the desired
vacuum subtraction.  Legacy files without sufficient provenance cannot be
automatically classified by their gamma channel alone.

The disconnected EMT derivative is not affected by this particular omission:
its derived quantity explicitly includes the minus in
Eq.~\eqref{eq:disc-emt-derivative}.  Its primitive local trace still follows
the positive-trace warning.

\section{Pion qTMD wave function and quasi-distribution amplitude}

\subsection{Vacuum-to-pion bilocal}

The sink bilocal realized by the qTMDWF/qDA contraction is
\begin{equation}
 \cO_\Gamma(x;b)
 =\bar d(x+b)W[x+b,x]\Gamma u(x).
 \label{eq:qTMDWF-operator}
\end{equation}
With a local pion source
\(\bar u(x_0)\Gsrc d(x_0)\), flavor-preserving Wick pairing gives
\begin{equation}
 C_\Gamma(x;b)
 \propto
 \Trsc\left[
 S_d(x_0,x+b)W[x+b,x]\Gamma
 S_u(x,x_0)\Gsrc
 \right],
 \label{eq:qTMDWF-wick}
\end{equation}
with the same pseudoscalar source-adjoint convention as
Eq.~\eqref{eq:pion-C2-production}.  Its spectral limit is
Eq.~\eqref{eq:vac-pion-spectral}.

\subsection{Why shift before dagger gives the correct endpoint}

Production first shifts or transports the source-to-sink backward-flavor
propagator:
\begin{equation}
 \widetilde S_d(x,x_0)
 =W[x,x+b]S_d(x+b,x_0).
 \label{eq:qTMDWF-pre-dagger}
\end{equation}
It then forms the backward line:
\begin{align}
 \gamma_5\widetilde S_d(x,x_0)^\dagger\gamma_5
 &=
 \gamma_5S_d(x+b,x_0)^\dagger
 W[x,x+b]^\dagger\gamma_5
 \nonumber\\
 &=S_d(x_0,x+b)W[x+b,x].
 \label{eq:qTMDWF-post-dagger}
\end{align}
The Wilson line is color-only and commutes with \(\gamma_5\).  Substitution
of Eq.~\eqref{eq:qTMDWF-post-dagger} into the common pion trace produces
Eq.~\eqref{eq:qTMDWF-wick}.  Thus the shift--dagger order does not reverse the
physical operator; it creates the required antiquark endpoint and conjugate
link.

\subsection{qDA as the straight bT=0 limit}

The gauge-invariant qDA restricts \(b_T=0\) and uses the straight transporter
of Eqs.~\eqref{eq:straight-positive}--\eqref{eq:straight-negative}.  A common
leading-twist continuum definition is
\begin{align}
 \avg{0|\bar d(z/2)\gamma_z\gamma_5
 W[z/2,-z/2]u(-z/2)|\pi(P)}
 &=
 \ii f_\pi P_z
 \int_0^1\dd x\,
 \ee^{\ii(x-\frac12)zP_z}
 \phi_\pi(x,\mu)
 \nonumber\\
 &\quad+\text{higher twist},
 \label{eq:pion-DA}
\end{align}
up to Euclidean continuation and normalization conventions.  Therefore the
axial channels, such as \code{Z5} or \code{T5} for a specified quasi-DA
choice, carry the leading-twist DA interpretation.  A scan of all 16 gamma
matrices is useful operator data but is not sixteen independent
leading-twist DAs.

The GI qDA endpoint and straight path are correct.  The CG qDA and the current
qTMDWF implementation use ordinary shifts and require fixed-gauge input.
Both remain bare vacuum-to-pion matrix elements until nonlocal
renormalization and matching are performed.  A gauge-invariant qTMDWF would
also require a staple and soft/rapidity scheme compatible with the target
factorization theorem.

\section{Legacy pion soft-factor four-point measurement}

\subsection{Wall-source propagators}

For a Coulomb-gauge wall source at \(t_0\) and quark momentum
\(\boldsymbol k\),
\begin{equation}
 \eta_{\boldsymbol k}(\boldsymbol x,t)
 =\delta_{t,t_0}\ee^{+\ii\boldsymbol k\cdot\boldsymbol x},
\qquad
 G_{\boldsymbol k}(x;t_0)=D^{-1}\eta_{\boldsymbol k}.
 \label{eq:wall-propagator}
\end{equation}
The legacy contraction uses both \(+\boldsymbol k\) and
\(-\boldsymbol k\) wall propagators and constructs backward antiquark lines
with gamma-five Hermiticity.

\subsection{Four-propagator Wick blocks}

For source time \(t_0\) and sink separation \(t_{\rm sep}\), define
\begin{align}
 G_w&=G_{k_{\rm fw}}(t_0),&
 G_{w,b^\dagger}&=G_{-k_{\rm bw}}(t_0),\\
 G_{w,\dagger}&=G_{-k_{\rm fw}}(t_0+t_{\rm sep}),&
 G_{w,b}&=G_{k_{\rm bw}}(t_0+t_{\rm sep}).
\end{align}
After the legacy momentum-transfer phase on the sink-side backward
propagator and an ordinary transverse shift by \(b\), the two closed
spin--color blocks are
\begin{align}
 A_b(x)
 &=G_w(x)\Gamma_{\rm src}\gamma_5
   G_{w,b^\dagger}(x+b)^\dagger\gamma_5,
 \label{eq:soft-A}\\
 B_b(x)
 &=G_{w,b}(x+b)\Gamma_{\rm sink}\gamma_5
   G_{w,\dagger}(x)^\dagger\gamma_5.
 \label{eq:soft-B}
\end{align}
The measured four-point correlator is
\begin{equation}
 C_4(t;t_{\rm sep},b,\Gamma_1,\Gamma_2)
 =\sum_{\boldsymbol x}
 \Trsc\left[
 A_b(x)\Gamma_2B_b(x)\Gamma_1
 \right].
 \label{eq:soft-C4}
\end{equation}
The production \code{tmp_1}, \code{tmp_2}, and final
\code{einsum} implement Eqs.~\eqref{eq:soft-A}--\eqref{eq:soft-C4} in this
order.  The legacy pion four-propagator Wick algebra is therefore
internally consistent.

\subsection{What the four-point contraction does not prove}

Equation~\eqref{eq:soft-C4} is a pion wall-source four-point function with
ordinary fixed-gauge transverse shifts.  A universal TMD soft function is
normally a vacuum expectation value of Wilson-line geometry defined in the
same rapidity regulator and renormalization scheme as the unsubtracted TMD.
Correct contraction of Eq.~\eqref{eq:soft-C4} alone does not prove
\begin{equation}
 C_4\equiv S_{\rm vacuum}^{\rm scheme}
\end{equation}
or prove that dividing the production GI staple matrix element by a ratio
built from \(C_4\) performs the desired soft subtraction.  That identification
requires an explicit factorization derivation matching staple length,
directions, rapidity regulator, gauge convention, and renormalization scheme.

\section{Background-current and current--current response}

\subsection{First derivative of a propagator}

Let
\begin{equation}
 D(\lambda)=D+\lambda O,\qquad S(\lambda)=D(\lambda)^{-1}.
\end{equation}
Differentiating \(D(\lambda)S(\lambda)=1\) gives
\begin{equation}
 \left.\frac{\dd S}{\dd\lambda}\right|_0
 =-SOS.
 \label{eq:first-response}
\end{equation}
For a local current with momentum transfer,
\begin{equation}
 O_q(x)=\Gamma_{\rm current}\Phi_{\boldsymbol q}(x;x_0),
\end{equation}
the inserted source is
\begin{equation}
 \eta_O(x)=O_q(x)S(x,x_0),
\end{equation}
and its inversion yields
\begin{equation}
 D^{-1}\eta_O=SOS.
 \label{eq:SOS-positive}
\end{equation}
\code{invert_local_current_response_propagator} constructs the positive
quantity in Eq.~\eqref{eq:SOS-positive} by default.  It is appropriate for a
direct comparison to an explicitly summed insertion.  Interpreting it as the
derivative of \(D+\lambda O\) requires the minus in
Eq.~\eqref{eq:first-response}; the helper exposes this as
\code{response_sign}.

If only a source-relative time window \(\mathcal T\) is selected, the source
is
\begin{equation}
 \eta_O^{\mathcal T}
 =\sum_{\tau\in\mathcal T}
 \delta_{x_4,t_0+\tau}\,O_q(x)S(x,x_0).
\end{equation}
The one inversion then produces the insertion summed over that window.  It is
not an unsummed \(C_3(\tau)\).

\subsection{Ordered second response}

For
\(D(\lambda_1,\lambda_2)=D+\lambda_1O_1+\lambda_2O_2\),
\begin{equation}
 \left.
 \frac{\partial^2S}{\partial\lambda_2\partial\lambda_1}
 \right|_0
 =SO_2SO_1S+SO_1SO_2S.
 \label{eq:mixed-second-response}
\end{equation}
The nested helper constructs the single ordered term
\begin{equation}
 SO_2SO_1S.
 \label{eq:ordered-response}
\end{equation}
It is a current--current response diagnostic or one time ordering, not by
itself the full mixed derivative in
Eq.~\eqref{eq:mixed-second-response}.  A physical second-order analysis must
define the second ordering, time windows, coincident-current contact terms,
and the spectral decomposition independently of a direct three-point
formula.

\section{Audit evidence and formula-to-code verdicts}

\subsection{Independent CPU oracles}

The numerical oracles use explicit NumPy sums and do not import the production
contraction functions.  Their double-precision algebraic tolerance was
\(10^{-11}\).

\begin{longtable}{L{0.24\textwidth}L{0.48\textwidth}L{0.20\textwidth}}
\toprule
Oracle & Formula tested & Maximum discrepancy/result\\
\midrule
\endfirsthead
\toprule
Oracle & Formula tested & Maximum discrepancy/result\\
\midrule
\endhead
Pion Wick oracle &
Full gamma-five backward line and
\(C_2^\pi=\Tr(S^\dagger S)\) &
\(5.68\times10^{-14}\)\\
Proton Wick oracle &
Complete flavor-labelled direct--exchange determinant and
\(T_2^{\rm code}=-\mathcal X\) &
\(5.55\times10^{-14}\) for full \(C_2\)\\
Sequential oracle &
\(K_d=\partial C_2/\partial S_d\) and
\(K_u=\partial C_2/\partial S_u\) &
\(1.99\times10^{-15}\), \(2.08\times10^{-14}\)\\
Production adapter &
Actual production down/up kernels and pion backward line versus independent
reference &
\(3.55\times10^{-15}\), \(3.56\times10^{-15}\); pion exact\\
Disconnected oracle &
Complete-basis trace estimator, determinant loop sign, and covariance &
Trace error \(4.48\times10^{-16}\); physical covariance exactly opposite\\
Geometry oracle &
Positive/negative straight link, staple chain, plane-wave derivative,
Fourier phase, and SU(3) covariance &
\(3.38\times10^{-15}\)\\
\bottomrule
\end{longtable}

\subsection{Single-A100 S8T8 evidence}

On a single A100 with \code{mpi_geometry=1.1.1.1}, the audit collected 63
relevant tests: 61 passed and two were skipped only because optional
pre-generated comparison directories were absent.  A real HYP-smeared
\(S8T8\) gauge field confirmed the cached GI staple against stepwise
\code{covDev}, including positive and negative \(b_z\), transverse staples,
and the PDF limit, at \(10^{-12}\) tolerance.

The source and pion-sequential solves reached true residuals
\(4.77\times10^{-13}\) and \(7.47\times10^{-13}\), respectively.  Pion and
proton \(C_2\), pion EMFF at \(q=0,1\), GI qGTMD, GI qGPD, CG qTMDWF, GI qDA,
and the symmetric derivative all produced finite arrays.

A four-rank run was not required to resolve any remaining formula conflict:
the compound-index, phase, and Wilson-path questions were already decided by
explicit oracles and real-gauge single-rank tests.  This is not a claim of
single-/multi-rank bitwise equivalence; a future MPI boundary regression
should compare \code{1.1.1.1}, \code{1.1.2.2}, and \code{2.2.1.1} when that
parallel property is the target.

\subsection{Source-level map}

\begin{longtable}{L{0.30\textwidth}L{0.38\textwidth}L{0.24\textwidth}}
\toprule
Production location & Formula represented & Verdict\\
\midrule
\endfirsthead
\toprule
Production location & Formula represented & Verdict\\
\midrule
\endhead
\code{pion_utils_vibe_develop.py:154--157} &
\(S(x_0,x)=\gamma_5S(x,x_0)^\dagger\gamma_5\) with full endpoint reversal &
\good{Correct}\\
\code{pion_utils_vibe_develop.py:236--246} &
\(B\Gsink S\Gsrc\) spin--color trace and Fourier sum &
\good{Correct}\\
\code{proton_utils_vibe_develop.py:87--176} &
Proton direct--exchange determinant &
\good{Correct}; printed term 2 is \(-\mathcal X\)\\
\code{bw_seq_pyquda.py:143--176} &
Meson fixed-sink source
\(\Phi_{p_f}\gamma_5\Gsink^\dagger\gamma_5S\) &
\good{Correct}\\
\code{bw_seq_pyquda.py:181--300} &
Down/up proton sequential cuts &
\good{Correct}; two/four terms\\
\code{pion_EMFF_vibe_develop.py:172--190} &
One-active-line local-current pion \(C_3\) &
\warn{Correct raw correlator}, not \(F_\pi\)\\
\code{qtmd_operator_utils.py:102--222} &
Straight and fixed-length staple transport &
\good{Correct geometry}\\
\code{application/nucleon_TMD/shared_runner.py:230--260} &
Proton nonlocal insertion trace &
\good{Correct raw connected trace}\\
\code{pion_qTMDWF_pyquda.py:130--240} &
Shifted backward line for qDA/qTMDWF &
\good{Correct endpoint}; CG needs gauge fixing\\
\code{flowed_fermion_bilinear_vibe_develop.py:148--181} &
Central derivative and left-line reconstruction &
\good{Correct lattice derivative}\\
\code{Disconnected_1pt_EMT_vibe_develop.py:302--329} &
Two-sided derivative loop &
\good{Closed-loop sign included}\\
\code{Disconnected_1pt_qTMD_vibe_develop.py:85--160} &
Positive stochastic qTMD/PDF trace &
\bad{Physical loop needs an outer minus}\\
\code{Disconnected_1pt_EMT_vibe_develop.py:1025--1110} &
Clover \(F\) and \(2\Tr FF\) &
\warn{Correct primitive}, not complete EMT\\
\code{pion_soft_factor_vibe_develop.py:208--280} &
Legacy two-block pion \(C_4\) &
\warn{Correct specified \(C_4\)}; universal soft identification unproved\\
\code{pion_current_background_response_vibe_develop.py:49--175} &
\(SOS\) and ordered \(SO_2SO_1S\) &
\good{Correct diagnostics}; derivative sign must follow definition\\
\bottomrule
\end{longtable}

\section{Measurement-by-measurement conclusions and data impact}

\begin{longtable}{@{}L{0.17\textwidth}L{0.19\textwidth}L{0.34\textwidth}L{0.22\textwidth}@{}}
\toprule
Measurement & Verdict & Physics qualification & Existing-data consequence\\
\midrule
\endfirsthead
\toprule
Measurement & Verdict & Physics qualification & Existing-data consequence\\
\midrule
\endhead
Pion common \(C_2/C_3\) &
\good{Correct} &
Flavor path, gamma-five backward line, and spin--color trace agree. &
No Wick-based recomputation required.\\
Proton common \(C_2\) &
\good{Correct} &
Full \(u\)-line direct--exchange determinant; no \(u/d\) cross contraction. &
No recomputation from the exchange concern.\\
Proton up/down sequential &
\good{Correct} &
Exact flavor derivatives of the same \(C_2\); up multiplicity already included. &
Do not multiply the up channel by two again.\\
Pion EMFF &
\warn{Conditionally correct} &
Raw one-line local-current \(C_3\); missing complete physical extraction. &
Retain \(C_3\); add missing \(C_2\), charges, \(Z_V\), improvement, and ratio.\\
Proton EMFF &
\warn{Not a dedicated target} &
Common local insertion exists, but no full \(F_1/F_2\) workflow. &
Do not relabel generic local-limit output as \(F_1,F_2\).\\
Connected quark EMT &
\warn{Conditionally correct} &
Derivative and Wick primitives correct; flowed and unrenormalized. &
Retain primitives; perform trace basis, ringed normalization, matching, mixing.\\
Disconnected quark EMT &
\warn{Conditionally correct} &
Derived derivative includes the loop sign; primitive local trace does not. &
Use dataset-level convention metadata.\\
Gluon EMT &
\warn{Conditionally correct} &
\(2\Tr FF\) primitive only. &
Retain primitive; construct trace/matching basis in analysis.\\
Connected GI qTMD/qPDF &
\warn{Correct raw kernel} &
Wick and path correct; current forward runner has \(P_z=0\), with no soft/matching. &
Raw data valid, but not a LaMET distribution; large-\(P_z\) target requires new data.\\
Connected qGTMD/qGPD &
\warn{Correct raw kernel} &
\(q\ne0\) is off-forward with
\(P=p_f-q/2\), \(\Delta=q\). &
Retain raw data and label/store \(P,\Delta,\xi,t\).\\
Disconnected qTMD/PDF &
\bad{P0 sign issue} &
Saved trace is positive; physical Wick loop is negative. &
Current-schema raw loops can be exactly sign-corrected without new inversions.\\
qTMDWF &
\warn{Conditionally correct} &
Endpoint Wick path correct; current implementation is CG/fixed-gauge. &
Usable only with gauge-fixing provenance and later renormalization/matching.\\
qDA &
\warn{Conditionally correct} &
GI \(b_T=0\) straight endpoint correct; leading twist needs axial channel. &
Bare GI correlator can be retained for later analysis.\\
Pion soft \(C_4\) &
\warn{Specified algebra correct} &
No demonstrated equality to the GI quasi-TMD vacuum soft function. &
Retain as a legacy \(C_4\); do not divide into GI qTMD without factorization proof.\\
Background response &
\good{Correct diagnostic} &
Positive \(SOS\) versus derivative \(-SOS\) is explicit; nested helper is one ordering. &
Interpret using response-sign and time-window metadata.\\
\bottomrule
\end{longtable}

\section{Compact derivation checklist for new analyses}

For every output channel, the following checks should be completed in order:
\begin{enumerate}
  \item Write the continuum operator and hadron matrix-element decomposition.
        State flavor charges and external-state normalization.
  \item Decide whether the file contains a raw correlator, fitted bare matrix
        element, or renormalized/matched observable.
  \item Use the actual plus-sign \code{MomentumPhase} convention to derive
        \(p_i=p_f-q\).  Store \(P,\Delta,\xi,t\) for off-forward data.
  \item Preserve flavor labels in the Wick derivation.  Apply
        \(S_u=S_d\) only after forbidden cross-flavor contractions have been
        excluded.
  \item For a proton, verify the complete direct--exchange determinant and
        use a flavor derivative to define every sequential source.
  \item For a disconnected bilinear, distinguish the positive stochastic
        trace from the negative physical closed loop before vacuum
        subtraction.
  \item For a bilocal, prove the endpoint and Wilson-link orientation after
        every dagger; distinguish fixed-gauge CG from gauge-invariant GI.
  \item For flowed EMT data, keep trace subtraction, ringed normalization,
        small-flow matching, and quark--gluon mixing explicit.
  \item Only after these steps construct ratios, fits, renormalization,
        matching, soft subtraction, and continuum/flow-time limits.
\end{enumerate}

\section{Final physics statement}

Starting from physical interpolators and preserving field order and flavor,
the connected pion and proton production kernels implement the intended Wick
paths.  In particular:
\begin{equation}
 C_2^\pi=\Trsc(S^\dagger S)
\end{equation}
in the equal-line pseudoscalar limit,
\begin{equation}
 C_2^p=-\mathcal D+\mathcal X
\end{equation}
contains the full same-flavor \(u\) determinant, and
\begin{equation}
 K_f=\frac{\partial C_2^p}{\partial S_f}
\end{equation}
generates the correct proton sequential sources.  Straight and staple
transporters have the intended endpoints, and shifting the qTMDWF/qDA
backward-flavor line before gamma-five Hermitian conjugation yields the
correct conjugate Wilson line.

The one confirmed contraction-level physical problem is
\begin{equation}
 \avg{\bar q\Gamma\cW q}
 =-\Tr(\Gamma\cW D^{-1}),
\end{equation}
whereas the disconnected qTMD/PDF producer saves the positive stochastic
trace and its documented physical covariance omits the outer minus.  This is
an exact convention error, not a statistical or solver effect.

All other major qualifications in this note concern the distinction between
a correct raw lattice building block and a final observable: EMFF requires
charge/current normalization and a ratio or spectral fit; EMT requires the
full flowed matching and mixing program; and nonlocal parton observables
require suitable external momentum, renormalization, matching, and, for TMDs,
a compatible soft and rapidity treatment.

\begin{thebibliography}{99}

\bibitem{DetmoldOrginos}
W.~Detmold and K.~Orginos,
``Nuclear correlation functions in lattice QCD,''
\href{https://arxiv.org/abs/1207.1452}{arXiv:1207.1452}.

\bibitem{Suzuki}
H.~Suzuki,
``Energy--momentum tensor from the Yang--Mills gradient flow,''
\href{https://arxiv.org/abs/1304.0533}{arXiv:1304.0533}.

\bibitem{MakinoSuzuki}
H.~Makino and H.~Suzuki,
``Lattice energy--momentum tensor from the Yang--Mills gradient flow:
inclusion of fermion fields,''
\href{https://arxiv.org/abs/1403.4772}{arXiv:1403.4772}.

\bibitem{QuasiGPD}
S.~Bhattacharya et al.,
``Generalized parton distributions from lattice QCD with asymmetric
momentum transfer,''
\href{https://arxiv.org/abs/2209.05373}{arXiv:2209.05373}.

\bibitem{QuasiTMD}
P.~Shanahan, M.~Wagman, and Y.~Zhao,
``Collins--Soper kernel for TMD evolution from lattice QCD,''
\href{https://arxiv.org/abs/2107.11930}{arXiv:2107.11930}.

\end{thebibliography}

\end{document}
